% ==========================================================================
% Section 5: Implementation and Experimental Setup
% ==========================================================================
\section{Implementation and Experimental Setup}
\label{sec:implementation-setup}

This section states the controls and scoring rules needed to interpret the
ablations in Section~\ref{sec:ablation-study}. The experiment families share
a benchmark but are not a single factorial experiment: candidate pools,
context construction, and latency accounting differ. Variants are therefore
compared only within their own controlled family.

\subsection{Pipeline and Frozen Artifacts}
\label{sec:software-architecture}

The evaluated pipeline has five stages: query embedding with
\texttt{intfloat/multilingual-e5-large}; dense FAISS retrieval and, where
enabled, tokenized BM25 retrieval; RRF fusion, graph expansion, or
cross-encoder reranking; serialization of retained chunks into a grounded
prompt; and answer generation and evaluation. The analysis reads frozen
manifests, predictions, aggregate metrics, latency records, and error logs;
it does not rerun models. Every primary comparison is joined by
\texttt{qa\_id}, so an aggregate difference is traceable to paired cases.

\subsection{Indexes, Graph, and Runtime Settings}
\label{sec:build-process}

The dense/hybrid/graph experiment uses a FAISS index with 1,513,376 chunk
vectors and a graph over the same number of chunks. The hybrid branch searches
ten BM25 shards. Dense/hybrid retrieval uses a 50-result search budget, broad
filtering, a final context of 10 chunks, and RRF parameter $k=60$ where BM25
is enabled. Graph variants take five seeds, traverse one hop, and retain at
most 10 graph-context chunks. The dense/hybrid and reranker generation
artifacts use GPT-4o-mini.

The reranker family fixes a 30-chunk RRF candidate pool (C30) and a 10-chunk
final context (T10), comparing no reranking, RRF-only, mMiniLM, Qwen3, BGE,
and Jina. The primary generator study compares GPT-4o-mini Base and CoT with
the saved top-10 context fixed by question. The metadata study changes only
the indexed string from chunk body to header-plus-chunk text; encoder, FAISS,
top-$k=10$, generator, prompt, seed, corpus, and benchmark stay fixed.

\subsection{Benchmark and Metrics}
\label{sec:evaluation-metrics}

The benchmark contains 500 questions: 400 answerable questions with gold
chunk identifiers and 100 unanswerable questions. Retrieval scores use only
the 400 answerable records. Given gold chunks $G$ and retrieved top-$k$ chunks
$R_k$, context recall is $|G \cap R_k|/|G|$; Hit@k indicates that at least one
gold chunk occurs. MRR and nDCG reward early gold rank. Candidate Recall@30 is
the evidence ceiling for the reranker study and is intentionally identical
for every reranker.

Exact Match, Token F1, and ROUGE-L are lexical overlap with reference answers.
The saved unanswerable score is a refusal-marker diagnostic, not a semantic
assessment of legal appropriateness. Citation presence and structural coverage
measure saved marker/parser fields, not claim-to-source entailment. Latency is
reported as saved p50/p95 timing; cached-BM25 reranker timing is not full
online sparse-retrieval timing.

\subsection{Pairing, Uncertainty, and Claim Boundaries}
\label{sec:experimental-protocol}

Headline effects are paired by \texttt{qa\_id}. Retrieval uses the shared
answerable intersection; scheduled-denominator answerability diagnostics keep
all 500 cases. Percentile paired-bootstrap confidence intervals use 10,000
resamples for headline metadata, dense/hybrid, reranker, and GPT comparisons;
reranker slices use 2,000. Intervals crossing zero are reported as
inconclusive. The non-GPT prompt pairs are descriptive when code provenance,
timeouts, retries, or retrieved context differ.

No blinded legal-correctness labels, claim-level faithfulness labels, or
validated semantic judge are available. Thus, no result is interpreted as
proof of legal correctness, legal safety, citation entailment, or faithful
reasoning merely because it improves a lexical, rank, citation-format, or
refusal-format metric.
